\section{Introduction}

Token log-probabilities are free. Every autoregressive language model already computes them during generation, and a long tradition of work \citep{kadavath2022,fadeeva2024,farquhar2024} treats the aggregated sequence-level statistic as a confidence proxy for hallucination detection. Yet the aggregation step itself --- the choice of \emph{which} scalar to extract from the per-token log-probability sequence --- has never been the subject of systematic study. Most methods pick one function (mean in CCP \citep{fadeeva2024}, sum-within-clusters in Semantic Entropy \citep{farquhar2024}) without testing whether a different choice would perform better, and without a theoretical reason to expect it would matter.

It matters by 12 AUROC points.

We show that taking the \emph{minimum} token log-probability outperforms the \emph{mean} by up to 0.119 AUROC on factual recall benchmarks (TriviaQA), while the \emph{mean} outperforms the \emph{minimum} by up to 0.120 AUROC on imitative-falsehood benchmarks (TruthfulQA) --- with effect sizes confirmed across two model families via bootstrap 95\% confidence intervals. The direction of advantage flips cleanly between benchmark types, and the flip is not an empirical coincidence: it is predicted by the structural nature of the hallucination being detected.

\textbf{The mechanism is as follows.} Recall-failure hallucinations --- where a model fails to retrieve a specific fact --- produce \emph{peaked} token probability distributions. The model generates function words with high probability (``The capital of France is'') but concentrates uncertainty at the single incorrect fact-token. The minimum log-probability captures this worst-case signal directly. Imitative-falsehood hallucinations --- where a model confidently generates an answer it has learned from training data but that is factually wrong --- produce \emph{flat} high-probability distributions. No single token carries anomalous uncertainty; instead, the hallucination signal is distributed uniformly across all tokens. The mean log-probability integrates this distributed signal that the minimum cannot resolve. The aggregation function must match the distribution shape, and the distribution shape is determined by the hallucination type.

This mechanism is empirically grounded in three layers of evidence. First, token probability distributions on TriviaQA hallucinated responses are significantly more peaked than correct responses (peakedness ratio: 2.936 vs.\ 2.533, $p = 0.002$, $n = 488$). Second, rank-order correlation confirms the direction at the signal level: Spearman $\rho(\min) > \rho(\text{mean})$ on TriviaQA and $\rho(\text{mean}) > \rho(\min)$ on TruthfulQA for both LLaMA-2-7B and Mistral-7B-v0.1, with all four bootstrap 95\% CIs excluding zero. Third, AUROC threshold analysis confirms the pattern with effect sizes of 0.056--0.119 (factual recall) and 0.111--0.120 (imitative falsehood).

An unexpected finding strengthens the practical contribution: raw unnormalized log-probability sum achieves the highest AUROC on TriviaQA for both models (0.895--0.896), exceeding both min (0.849--0.892) and mean (0.730--0.836) --- and exceeding published Semantic Entropy AUROC ($\approx$0.79) \citep{farquhar2024} at 10$\times$ lower inference cost (under our evaluation protocol; see Section~6.3 for cross-pipeline caveats). Contrary to our initial prediction, the length-sensitive ``worst'' aggregator is the strongest factual-recall detector, suggesting sequence-level joint probability as an underexplored zero-cost hallucination signal.

\textbf{Contributions.} Building on the hallucination-type mechanism, this paper makes four contributions:

\begin{enumerate}
\item \textbf{Mechanistic taxonomy.} We identify hallucination type (recall-failure vs.\ imitative-falsehood) as the determinant of token probability distribution shape (peaked vs.\ flat), and distribution shape as the determinant of optimal aggregation function (min vs.\ mean). We provide the first controlled ablation isolating aggregation function as the sole experimental variable on fixed factual QA splits, filling the gap identified as open by \citet{liu2025}.

\item \textbf{Distributional evidence.} We provide direct experimental confirmation that recall-failure hallucinations produce significantly higher token distribution peakedness than correct responses at 7B model scale ($p = 0.002$), converting the mechanism's first step from a theoretical assumption to an empirical fact.

\item \textbf{Aggregation selection criterion.} We establish a principled, label-free selection criterion for zero-cost hallucination detection: use min log-probability for factual recall tasks; use mean log-probability for imitative-falsehood tasks. Both require only a single forward pass with frozen model weights.

\item \textbf{Unexpected finding.} Raw unnormalized log-probability achieves higher AUROC than both normalized aggregations on factual recall benchmarks, and surpasses published multi-sample baselines under our evaluation protocol, positioning sequence-level joint probability as a novel zero-cost detection signal warranting further investigation.
\end{enumerate}

The remainder of the paper is organized as follows. Section~2 positions our work within the hallucination detection literature. Section~3 describes our experimental methodology, the mechanism hypothesis, and the four-hypothesis verification design. Section~4 presents the experimental setup. Section~5 reports results across all four hypotheses. Section~6 discusses implications, limitations, and future work. Section~7 concludes.
